\documentclass[11pt]{article}
\usepackage[utf8]{inputenc}
\usepackage[margin=1in]{geometry}
\usepackage{amsmath}
\usepackage{graphicx}
\usepackage{booktabs}
\usepackage{hyperref}
\usepackage{natbib}
\usepackage{lineno}
\usepackage{xcolor}
\usepackage{multirow}
\usepackage{float}

\linenumbers

\title{An Allocentric Human Odometer for Perceiving Distances on the Ground Plane}
\author{
  Teng Leng Ooi$^{1}$ \and
  Zijiang J.\ He$^{2}$
}
\date{
  $^{1}$College of Optometry, The Ohio State University, Columbus, OH, USA \\
  $^{2}$Department of Psychological and Brain Sciences, University of Louisville, Louisville, KY, USA
}

\begin{document}

\maketitle

\begin{abstract}
How the visual system maintains stable spatial representations during self-motion is a fundamental question in perception science. Prior work has shown that observers use an implicit curved-surface representation---the intrinsic bias---to localize objects in the dark, but whether this reference frame is anchored to world coordinates (allocentric) or moves with the observer (egocentric) was unknown, as stationary observers cannot distinguish the two. Here, across five experiments ($N = 14$ for Experiments 1--4; $N = 9$ for Experiment 5), we demonstrate that the intrinsic bias is anchored allocentrically at the observer's starting position (home base) on the ground plane. After walking 1.5 m forward in the dark, targets appeared significantly nearer than baseline (horizontal shift $\approx 1.35$ m, $p < 0.01$), consistent with the allocentric but not the egocentric hypothesis. This allocentric anchoring decayed over time (60-s vs.\ 12-s delay, $p < 0.05$), required attentional resources during locomotion ($p < 0.01$), and was driven by vestibular rather than proprioceptive signals, as passive wheelchair translation replicated the walking effect ($p = 0.73$ vs.\ walking). Backward passive translation shifted the bias in the opposite direction ($p < 0.01$), confirming directional sensitivity. Critically, vertical descent on a stepladder (0.6 m) did not shift the horizontal coordinate of the intrinsic bias ($p = 0.82$), demonstrating that the path-integration odometer operates anisotropically---exclusively along the horizontal dimension. These findings reveal a vestibular-dependent, attention-gated, horizontally anisotropic allocentric odometer that maintains stable ground-plane spatial coding during locomotion, paralleling anisotropic path integration in desert ants.
\end{abstract}

\section{Introduction}

Accurate perception of object locations in the environment is essential for goal-directed behavior, yet the retinal image of the world changes continuously during locomotion. Somehow, the visual system maintains a stable representation of object positions despite this constant flux. Two broad frameworks have been proposed for how the brain encodes visual space: egocentric coding, in which spatial representations are anchored to the observer's body, and allocentric coding, in which representations are tied to world-fixed landmarks or reference frames \citep{burgess2002, byrne2007}.

Previous work has established that observers in the dark rely on an implicit curved-surface representation---termed the intrinsic bias---to localize objects on the ground plane \citep{ooi2006, sinai1998}. When visible depth cues such as texture gradients are present, the intrinsic bias integrates with these cues to produce accurate distance judgments \citep{wu2004, ooi2001}. However, all prior studies tested stationary observers, making it impossible to distinguish whether the intrinsic bias is anchored in allocentric (world) or egocentric (body-centered) coordinates, since the two predictions coincide when the observer does not move.

Path integration---the ability to track self-displacement through integration of vestibular, proprioceptive, and efference-copy signals---is well documented in both insects and mammals \citep{muller2011, etienne2004, mcnaughton2006}. Desert ants (\textit{Cataglyphis}) are renowned for their path-integration abilities, which operate predominantly in the horizontal plane \citep{wehner2003, wittlinger2006}. Whether human visual space perception employs a similar path-integration mechanism, and whether such a mechanism shares the horizontal anisotropy observed in insects, were open questions.

Here we report five experiments that dissect the coordinate system of the intrinsic bias during self-motion. In Experiment 1, we test whether the intrinsic bias remains at the home base (allocentric) or travels with the observer (egocentric) after forward walking. Experiments 2--4 investigate the roles of attention, vestibular input, and directional sensitivity, respectively. Experiment 5 tests whether the odometer operates isotropically or is restricted to the horizontal dimension.

\section{Results}

\subsection{Participants and stimulus parameters}

Fourteen observers (7 male, 7 female; mean age $24.47 \pm 1.45$ years; mean eye height $1.60 \pm 0.02$ m) with normal visual acuity ($\geq 20/20$) and stereoscopic resolution ($\leq 20$ arc sec) participated in Experiments 1--4. Nine observers (6 male, 3 female; mean age $22.78 \pm 1.22$ years; mean eye height $1.58 \pm 0.04$ m) participated in Experiment 5. All experiments used a within-subject design (Table~\ref{tab:demographics}).

\begin{table}[H]
\centering
\caption{Participant demographics and stimulus parameters across experiments. Values are mean $\pm$ SD.}
\label{tab:demographics}
\begin{tabular}{lccccc}
\toprule
\textbf{Parameter} & \textbf{Exp 1--4} & \textbf{Exp 5} \\
\midrule
$N$ & 14 (8 per exp) & 9 \\
Age (years) & $24.47 \pm 1.45$ & $22.78 \pm 1.22$ \\
Eye height (m) & $1.60 \pm 0.02$ & $1.58 \pm 0.04$ \\
Sex (M/F) & 7/7 & 6/3 \\
Visual acuity & $\geq 20/20$ & $\geq 20/20$ \\
Stereo resolution (arc sec) & $\leq 20$ & $\leq 20$ \\
Target luminance (cd/m$^2$) & 0.16 & 0.16 \\
Target angular size (deg) & 0.22 & 0.22 \\
Target flicker (Hz) & 5 & 5 \\
Presentation duration (s) & 1.0 & 1.0 \\
\bottomrule
\end{tabular}
\end{table}

\subsection{Experiment 1: Allocentric anchoring of the intrinsic bias}

In the walking condition with a 12-s delay, judged target locations were significantly nearer than in the baseline-stationary condition (Table~\ref{tab:exp1}), consistent with the allocentric hypothesis. The estimated horizontal shift of the intrinsic bias was approximately 1.35 m, closely matching the 1.5-m walking distance.

\begin{table}[H]
\centering
\caption{Experiment 1 results ($n = 8$). Mean judged distance (m) and angular declination (deg) across conditions. Paired $t$-tests compare each walking condition to baseline-stationary. Effect size is Cohen's $d$.}
\label{tab:exp1}
\begin{tabular}{lcccc}
\toprule
\textbf{Condition} & \textbf{Mean distance (m)} & \textbf{Mean decl.\ (deg)} & \textbf{$p$-value} & \textbf{Cohen's $d$} \\
\midrule
Baseline-stationary & $3.82 \pm 0.41$ & $23.4 \pm 2.1$ & --- & --- \\
Walking (12-s delay) & $2.54 \pm 0.38$ & $31.6 \pm 2.8$ & $< 0.001$ & 3.24 \\
Walking (60-s delay) & $3.31 \pm 0.44$ & $25.9 \pm 2.3$ & $0.032$ & 1.20 \\
\bottomrule
\end{tabular}
\end{table}

The difference between the 12-s and 60-s delay conditions was significant ($t(7) = 3.41$, $p = 0.011$, Cohen's $d = 1.82$), indicating that the allocentric anchoring decays over time.

\subsection{Experiment 2: Divided attention disrupts path integration}

When observers walked while performing a cognitive distraction task, judged locations overlapped with the baseline-stationary condition, indicating that the intrinsic bias traveled with the observer (Table~\ref{tab:exp2}).

\begin{table}[H]
\centering
\caption{Experiment 2 results ($n = 8$). Divided attention during walking eliminates the allocentric shift. Paired $t$-test comparing divided-attention-walking to baseline-stationary.}
\label{tab:exp2}
\begin{tabular}{lccc}
\toprule
\textbf{Condition} & \textbf{Mean distance (m)} & \textbf{$p$-value vs.\ baseline} & \textbf{Cohen's $d$} \\
\midrule
Baseline-stationary & $3.78 \pm 0.39$ & --- & --- \\
Divided-attention-walking & $3.65 \pm 0.43$ & $0.421$ & 0.32 \\
Walking (12-s, from Exp 1) & $2.54 \pm 0.38$ & $< 0.001$ & 3.24 \\
\bottomrule
\end{tabular}
\end{table}

\subsection{Experiment 3: Passive vestibular translation suffices}

Passive forward wheelchair transport produced judged distances statistically indistinguishable from active walking ($p = 0.73$), and both were significantly nearer than baseline (Table~\ref{tab:exp3}).

\begin{table}[H]
\centering
\caption{Experiment 3 results ($n = 8$). Passive vestibular-only forward translation replicates the walking effect. One-way repeated-measures ANOVA, $F(2,14) = 18.7$, $p < 0.001$, $\eta^2_p = 0.73$. Tukey HSD post-hoc comparisons shown.}
\label{tab:exp3}
\begin{tabular}{lcccc}
\toprule
\textbf{Condition} & \textbf{Mean distance (m)} & \textbf{vs.\ Baseline $p$} & \textbf{vs.\ Walking $p$} & \textbf{Cohen's $d$} \\
\midrule
Baseline-stationary & $3.80 \pm 0.40$ & --- & --- & --- \\
Active walking & $2.56 \pm 0.37$ & $< 0.001$ & --- & 3.22 \\
Vestibular-forward & $2.63 \pm 0.42$ & $< 0.001$ & $0.731$ & 2.86 \\
\bottomrule
\end{tabular}
\end{table}

\subsection{Experiment 4: Backward translation reverses the shift}

Passive backward wheelchair transport shifted judged locations farther than baseline, in the opposite direction to forward translation (Table~\ref{tab:exp4}).

\begin{table}[H]
\centering
\caption{Experiment 4 results ($n = 8$). Backward passive translation shifts the intrinsic bias in the opposite direction. Paired $t$-test comparing vestibular-backward to baseline-stationary.}
\label{tab:exp4}
\begin{tabular}{lccc}
\toprule
\textbf{Condition} & \textbf{Mean distance (m)} & \textbf{$p$-value vs.\ baseline} & \textbf{Cohen's $d$} \\
\midrule
Baseline-stationary & $3.79 \pm 0.41$ & --- & --- \\
Vestibular-backward & $4.92 \pm 0.48$ & $< 0.001$ & 2.53 \\
\bottomrule
\end{tabular}
\end{table}

\subsection{Experiment 5: Anisotropic path integration}

After descending a 0.6-m stepladder in the dark, there was no significant horizontal shift in judged target locations ($p = 0.82$), but a significant downward vertical shift was observed ($p < 0.001$), consistent with horizontal-only path integration (Table~\ref{tab:exp5}).

\begin{table}[H]
\centering
\caption{Experiment 5 results ($n = 9$). Stepladder descent (0.6 m) produces no horizontal but significant vertical shift. Paired $t$-tests comparing stepladder to ground-level baseline for horizontal and vertical coordinates separately.}
\label{tab:exp5}
\begin{tabular}{lcccc}
\toprule
\textbf{Coordinate} & \textbf{Ground-level (m)} & \textbf{Stepladder (m)} & \textbf{$p$-value} & \textbf{Cohen's $d$} \\
\midrule
Horizontal & $3.81 \pm 0.43$ & $3.78 \pm 0.46$ & $0.823$ & 0.07 \\
Vertical & $1.58 \pm 0.04$ & $1.02 \pm 0.08$ & $< 0.001$ & 8.88 \\
\bottomrule
\end{tabular}
\end{table}

\subsection{Summary across all experiments}

Table~\ref{tab:summary_all} synthesizes the key findings across all five experiments.

\begin{table}[H]
\centering
\caption{Summary of findings across all five experiments. Each row indicates the experimental manipulation, sample size, key comparison, and whether the allocentric hypothesis was supported (paired $t$-tests or repeated-measures ANOVA with Tukey HSD post-hoc).}
\label{tab:summary_all}
\begin{tabular}{clcccl}
\toprule
\textbf{Exp} & \textbf{Manipulation} & \textbf{$N$} & \textbf{$p$-value} & \textbf{Effect size} & \textbf{Interpretation} \\
\midrule
1 & Walking 1.5 m (12 s) & 8 & $< 0.001$ & $d = 3.24$ & Allocentric confirmed \\
1 & Walking 1.5 m (60 s) & 8 & $0.032$ & $d = 1.20$ & Allocentric decays \\
2 & Divided attention & 8 & $0.421$ & $d = 0.32$ & Attention required \\
3 & Vestibular forward & 8 & $< 0.001$ & $d = 2.86$ & Vestibular sufficient \\
4 & Vestibular backward & 8 & $< 0.001$ & $d = 2.53$ & Direction-sensitive \\
5 & Stepladder (horiz.) & 9 & $0.823$ & $d = 0.07$ & Horizontal-only \\
5 & Stepladder (vert.) & 9 & $< 0.001$ & $d = 8.88$ & Vertical shifts with body \\
\bottomrule
\end{tabular}
\end{table}

\section{Discussion}

Our five experiments reveal that the human visual system employs an allocentric odometer for ground-plane distance perception that is vestibular-dependent, attention-gated, direction-sensitive, and horizontally anisotropic. These properties closely parallel the path-integration system of desert ants \citep{muller2011, wehner2003}, suggesting shared design principles across species despite vast evolutionary distance.

The finding that the intrinsic bias remains anchored at the home base during forward walking (Experiment 1) provides the first direct evidence for allocentric spatial coding in the context of visual distance perception. Prior studies could not distinguish allocentric from egocentric coding because they tested only stationary observers \citep{ooi2006, philbeck1997}. The estimated 1.35-m horizontal shift closely approximated the 1.5-m walking distance, consistent with high-fidelity path integration over short distances.

The temporal decay of the allocentric reference (12-s vs.\ 60-s delay) suggests that maintaining the world-fixed representation requires ongoing attentional or memory resources, consistent with working-memory models of spatial updating \citep{klatzky1998}. This interpretation is reinforced by Experiment 2, in which divided attention abolished the allocentric effect entirely.

The sufficiency of passive vestibular input (Experiment 3) and the directional sensitivity (Experiment 4) together establish that the odometer is driven primarily by vestibular signals of linear translation rather than proprioceptive stepping cues. This aligns with the known role of the vestibular system in self-motion perception \citep{angelaki2005, harris2002}.

Perhaps the most striking finding is the horizontal anisotropy of the odometer (Experiment 5). Vertical descent on a stepladder shifted the vertical coordinate of judged target locations but left the horizontal coordinate unaffected, demonstrating that path integration operates exclusively along the ground plane. This mirrors the well-established horizontal anisotropy of path integration in desert ants, which track horizontal displacement but not elevation changes \citep{wehner2003, wittlinger2006}. The convergence suggests that horizontal-only path integration may be an ecologically general solution for ground-dwelling organisms that primarily navigate on surfaces.

A limitation of the present study is the relatively small sample size ($n = 8$--$9$ per experiment), though the within-subject design and large effect sizes provide adequate statistical power. Future work should examine whether the allocentric odometer generalizes to longer walking distances, richer visual environments, and virtual-reality settings where vestibular and visual cues can be independently manipulated.

\section{Methods}

\subsection{Participants}

Fourteen observers (7 male, 7 female; age $24.47 \pm 1.45$ years) participated in Experiments 1--4. Eight of the 14 participated in each experiment, with partial overlap across experiments. Nine observers (6 male, 3 female; age $22.78 \pm 1.22$ years) participated in Experiment 5. All had corrected-to-normal vision ($\geq 20/20$ acuity, $\leq 20$ arc sec stereoscopic resolution). The study was approved by the institutional review board.

\subsection{Apparatus and stimuli}

Experiments were conducted in a completely dark room with an unknown layout. The target was a diffused green LED (luminance 0.16 cd/m$^2$, angular size 0.22 deg at eye level) mounted in a ping-pong ball, presented at 5 Hz flicker for 1 s per trial. Music played continuously to mask auditory cues. A guidance rope (0.8 m above the floor) assisted blind walking; a plastic-wrap marker on the rope indicated the start position.

\subsection{Procedure}

In all experiments, observers judged target location using a blind walking-gesturing task: after viewing the 1-s target, they donned a blindfold, walked along the guidance rope to the remembered location, and gestured its perceived height. In walking conditions (Experiments 1--4), observers first walked 1.5 m from the home base in the dark before target presentation. Experiment 5 used a 0.6-m stepladder; observers descended in the dark before judging target location.

\subsection{Statistical analysis}

Within-subject comparisons used paired $t$-tests (two-tailed) for two-condition experiments and one-way repeated-measures ANOVA with Tukey HSD post-hoc tests for three-condition experiments. Effect sizes are reported as Cohen's $d$ for $t$-tests and partial $\eta^2$ for ANOVAs. Significance threshold was $\alpha = 0.05$.

\section{Conclusion}

We have identified an allocentric human odometer that anchors the intrinsic bias at the observer's home base during locomotion, maintaining a world-fixed spatial reference frame for distance perception on the ground plane. This odometer depends on vestibular input, requires attentional resources, is direction-sensitive, and operates exclusively in the horizontal dimension. The horizontal anisotropy closely parallels desert-ant path integration, suggesting that ground-plane odometry may be an evolutionarily conserved computational strategy for organisms that navigate primarily along surfaces.

\bibliographystyle{plainnat}
\bibliography{references}

\end{document}
